% =====================================================================
\section{Limitations and conclusion}
\label{sec:end}

\subsection{Limitations}
\label{sec:limits}
\begin{itemize}
\item \textbf{Selection.} The evidence map is curated. The original searches were not logged query by query, and counts describe this map, not the prevalence of opportunities.
\item \textbf{No new experiments.} We performed no new physical simulations or hardware benchmarks. Reported values are reconciled and scenario arithmetic is provided; no scenario is a benchmark result.
\item \textbf{Depth of reading.} Primary-source extraction now covers the OTOC acquisition time, the Navier--Stokes Toffoli and rotation counts, two gauge-theory evolution counts and the plaquette-oracle counts, and the per-instance records of the drag study. These checks establish what the sources report, not that their assumptions hold. The drag study's published range and fit are not reconciled with its public per-instance files, whose tolerance is recorded only in a task string, and we have not yet asked the source's authors. The spectral-amplification schedule is still a second-hand value, because the article of record could not be read by our tools. Values known only from an abstract or at second hand are marked in the ledger and kept out of figures and ratios.
\item \textbf{Incomplete comparison records.} For most device and resource-estimate entries, the tolerance, the meaning of the failure probability and the access model have not been extracted from the sources; Appendix~\ref{app:ledger} gives the counts. Boson-sampling experiments are now recorded separately; later random-circuit experiments are not. Of the seven pair records, the two lattice-model pairs and the spectroscopy pair still lack the source's failure allocation.
\item \textbf{Classical baselines not costed.} No Monte Carlo time-to-solution exists for the half-filled Hubbard cylinder; we found no neural-state or GPU DMRG result on the lattices of the published crossover; and the match between the GPU DMRG iron--sulfur calculation and the hardware experiment is not established.
\item \textbf{A moving target.} Status is recorded as of 11 November 2026. Sources are revised, and a summary can differ from the table it summarizes: one source used here changed its gate and cycle counts between versions, and its abstract page and its PDF named different gate units. The ledger therefore identifies the version and the artifact read for every plotted value.
\item \textbf{One author.} The assessments are one author's reading of the literature, from the institutional position described under Competing interests.
\end{itemize}

\subsection{Conclusion}
\label{sec:conclusion}
The evidence reviewed here supports neither a general claim that AI makes quantum computing unnecessary nor a general claim that quantum hardware already outperforms the best classical workflows for scientific applications. Of \wTimeline{} device claims, \wTLvalidation{} received classical responses at their validation benchmark and \wTLchallenged{} were challenged on it, \wTLobservables{} at its reported observables and \wTLpartial{} in part; \wTLcontested{} are contested and \wTLopen{} remain unresolved in the regimes tested. The responses came from structured numerical and variational methods. Neural methods contribute by other routes, as direct wavefunction solvers and through reusable predictions and validated reformulation. Classical and neural methods continue to improve, and fault-tolerant quantum estimates have also fallen substantially for selected models, to hours on about a million physical qubits for two flagship molecules.

The practical question is which method can deliver a specified scientific answer at an acceptable cost and reliability. That comparison depends on the instance, the observable, the accuracy, the access to data, state preparation, repetitions and the surrounding workflow. The strongest current opportunities, in coherent dynamics, correlated lattice models, multireference chemistry, core-level spectroscopy and gauge theory, are therefore conditional research targets and not established application-wide rankings. For one large class the evidence points the other way: sparse linear systems and fluid dynamics at engineering scale, where complexity limits allow at most a polynomial speedup and published end-to-end estimates are far from practical. Query-complexity results explain why access assumptions and promises must be stated; they do not determine the runtime ranking of the workloads reviewed here.

The evidence ledger records the conditions behind each assessment and the tests that would change it. Progress will be most informative when both sides publish complete cost accounting, strong baselines and ablations that isolate their contribution to the scientific result.

% =====================================================================
\section*{Data and code availability}
\addcontentsline{toc}{section}{Data and code availability}
The evidence ledger and its scripts are provided as ancillary files with this submission: eight tables (\texttt{entries}, \texttt{pairs}, \texttt{resources}, \texttt{values}, \texttt{models}, \texttt{forward\_classes}, \texttt{figure\_points} and \texttt{update\_log}, each a CSV file); three scripts (\texttt{build\_ledger.py}, the authoritative source of the tables; \texttt{make\_figdata.py}, validation and generation; \texttt{test\_negative.py}, fourteen corruptions the validator must reject); and \texttt{README.md} and \texttt{CHANGELOG.md}. Running the first two scripts regenerates the plotted data and the numbers quoted in the text (Appendix~\ref{app:ledger}). The scripts check internal consistency and source locators. They do not validate the scientific judgments, and they support the reproducibility of this survey, not the correctness of any cited result.

\section*{Competing interests}
\addcontentsline{toc}{section}{Competing interests}
The author directs RIKEN R-CCS, which operates the supercomputer Fugaku, and leads the FugakuNEXT project. R-CCS participates in the RIKEN--IBM quantum-centric supercomputing programme. Three lines of work discussed here have RIKEN co-authors and used Fugaku: the sample-based diagonalization experiments and their closed-loop integration, the embedded protein--ligand workflow, and the Heron Floquet experiment, whose classical comparisons include Fugaku runs (Section~\ref{sec:devices}). The Quantum Advantage Tracker cited in several cases is described in~\cite{legeza2026} as maintained by IBM and RIKEN. The author had no technical role in any of these studies or in the tracker; his connection to them is institutional. Section~\ref{sec:strategy} is the author's perspective.

\section*{Use of AI tools}
\addcontentsline{toc}{section}{Use of AI tools}
Large language models were used extensively throughout this work. Models from Anthropic (Claude) were used for literature retrieval, reading and checking of sources, drafting and revision of the text, and generation of the figures, the evidence ledger and its scripts. Models from OpenAI (ChatGPT) were used for several rounds of adversarial review of the manuscript, ledger and scripts, and the revisions respond to those reviews. One result discussed in Section~\ref{sec:ai} was itself obtained by others with a non-public Anthropic model. The author directed the work, made the judgments recorded here, reviewed the content, and is solely accountable for its accuracy and conclusions; errors that remain are the author's responsibility, not that of any tool.
